% !TeX root = ../main.tex

\section{Investigating a Single Neuron}

\subsection{A Continuously Spiking Neuron}

As a first exercise, set the neuron up such that it spikes continuously without any external input. For that purpose, play
with the function’s parameters and note down a suitable configuration in the following cell. Make sure to also save the
resulting plot.



\subsection{Simple Spike Train Statistics}

Calculate the neuron’s average firing rate.
Derive the inter-spike intervals (ISIs) of the spike train recorded in the previous exercise. The inter-spike interval
denotes the time between two consecutive spikes of a neuron. Plot a histogram of the ISIs.
Identify a method to calculate the mean and standard deviation of the neuron’s instantaneous firing rate. The
instantaneous firing rate provides a moment-by-moment measure of the neuron’s activity.
Hint: You may use, e.g., np.diff, np.mean, and np.std to complete these tasks.



\subsection{Recording the f-I curve of a silicon neuron}

Next, we want to study the behavior of a single neuron stimulated with a current source. For this, we introduce a new
parameter for our neuron model. This will allow to enable/disable a constant current source and observe the impact.
Add \texttt{constant_current_enable} (True/False) and \texttt{constant_current_i_offset} \( (range=0-1022) \) as parameters into your experiment
Configure a non firing neuron
Vary the current and observe the impact
To quantify your observations write a program that sweeps over a current range from 0 up to 800.
Plot the neuron’s mean instantaneous firing rate including its error over current.


\subsection{Synaptic Stimuli and PSP Stacking}

Your goal is to reproduce the image of PSP-Stacking
(summation) from the introduction. For this, you might find it helpful to look into \texttt{pynn_introduction} again.
Set up a population of one neuron with the default neuron parameters to record its membrane potential.
Create a population of one SpikeSourceArray as a stimulating population.
Connect the stimulating population to the neuron population.
Find a suitable spike pattern to recreate the plot (you might require multiple spikes).
Can the neuron spike with the current neuron configuration? Explain.
Suggest and implement changes in the configuration so that the neuron spikes. Try to apply one change at a time.
Explain your choices.
Can you tell which changes can be theoretically equivalent?
Hint: Changes can be related to neuron parameters, projection, or population.


